\documentclass[12pt,a4paper]{article}
\usepackage{fontspec}
\usepackage{unicode-math}
\setmainfont{Linux Libertine O}
\setmathfont{TeX Gyre Pagella Math}
\setmonofont[Scale=0.85]{Everson Mono}
\usepackage[top=1in, bottom=1in, left=1.1in, right=1.1in, headheight=15pt]{geometry}
\usepackage{microtype}
\usepackage{parskip}
\usepackage[english]{babel}
\setlength{\emergencystretch}{2em}
\usepackage{fancyhdr}
\pagestyle{fancy}\fancyhf{}\fancyhead[L]{\leftmark}\fancyhead[R]{\thepage}
\renewcommand{\headrulewidth}{0.4pt}
\usepackage{titlesec}
\titleformat{\section}{\Large\bfseries}{\thesection}{1em}{}
\titleformat{\subsection}{\large\bfseries}{\thesubsection}{1em}{}
\usepackage{enumitem}
\setlist[itemize]{topsep=0.5ex, itemsep=0.25ex, parsep=0pt}
\setlist[enumerate]{topsep=0.5ex, itemsep=0.25ex, parsep=0pt}
\usepackage{amsmath,amsthm}
\usepackage{mathtools}
\usepackage{booktabs}
\usepackage{array}
\usepackage{tabularx}
\usepackage[table]{xcolor}
\renewcommand{\arraystretch}{1.3}
\usepackage{tcolorbox}
\usepackage{hyperref}
\usepackage{cleveref}
\hypersetup{colorlinks=true, linkcolor=blue, urlcolor=cyan, citecolor=blue}
\providecommand{\Box}{}\renewcommand{\Box}{\mathord{\text{\fontspec{Everson Mono}◻}}}
\newtheorem{theorem}{Theorem}[section]
\newtheorem{lemma}[theorem]{Lemma}
\newtheorem{proposition}[theorem]{Proposition}
\newtheorem{corollary}[theorem]{Corollary}
\newtheorem{definition}[theorem]{Definition}
\newtheorem{remark}[theorem]{Remark}

\begin{document}
\title{\textbf{The Fold Verdict of a Coding Sequence}\\[0.3em]
\large A closed form for a four-valued structural readout, and its reduction to two codon families}
\author{C.\ Lando Mills\\
Trabuco Canyon, California\\
\texttt{c.landonmills@gmail.com}\\
ORCID: \href{https://orcid.org/0000-0003-0003-0552}{0000-0003-0003-0552}}
\date{16 August 2026}
\maketitle

\begin{abstract}
A protein sequence can be read as a word in a twelve-letter alphabet. The reading
is not an analogy imposed from outside: the standard genetic code promotes exactly
twelve of the twenty amino acids to a distinguished set, that set is in bijection
with twelve structural operations, and the remaining eight activate none of them
and are silent. Two of the twelve operations are a three-way \emph{fork} and a
three-way \emph{fuse}; the other ten are boundaries, self-reference, and work.

Over such words there is a four-valued structural readout, taking values in the
Belnap lattice $\{T, B, N, F\}$: the word \emph{closes} ($T$) when every fork pairs
with a fuse around the cycle and some paired region contains work; it is
\emph{overdetermined} ($B$) when a fork is left open; \emph{underdetermined} ($N$)
when forks and fuses pair but enclose nothing that does work; and \emph{ill-typed}
($F$) when a fuse has no fork to pair with. The definition is given by a pairing
procedure, and it is not obvious from the definition what a word must look like to
receive each value.

The central distinction this paper establishes is that the readout depends on the
word through only two numbers and one adjacency condition. Let $|w|_{\in}$ and
$|w|_{\ni}$ count forks and fuses. The sign of $|w|_{\in} - |w|_{\ni}$ decides
between $B$ and $F$ outright, with the magnitude carrying no information; and when
that difference is zero, the word closes precisely when some fork is separated from
the fuse it pairs with by at least one of the seven letters that do work. Nothing
else about the arrangement matters for the value.

Carried back through the genetic code, the fork is asparagine and the fuse is
glutamine, so the readout of a coding sequence is a statement about two codon
families, $\mathrm{AAY}$ and $\mathrm{CAR}$, which differ in their first base
alone. We give the closed form, prove it from the pairing definition, exhibit the
witness that separates the correct enclosure condition from a weaker one that
agrees with it on every sequence we tested, and verify the result on a corpus of
sixty-five deposited structures and on coding sequences reconstructed from them.
\end{abstract}

\section{The alphabet and the four values}

\subsection{Twelve operations}

The alphabet is twelve marks. Each is an operation on a state, and the division
that matters here is between the marks that transform a state and the marks that do
not.

\begin{center}
\begin{tabular}{clll}
\toprule
Mark & Operation & Axis & Class \\
\midrule
$\vdash$   & source boundary        & Dimensionality  & no work \\
$\dashv$   & sink boundary          & Topology        & no work \\
$\odot$    & identity, self-reference & Criticality   & no work \\
$\in$      & three-way fork         & Granularity     & no work \\
$\ni$      & three-way fuse         & Grammar         & no work \\
\midrule
$\succ$    & forward morphism       & Recognition     & work \\
$\prec$    & reverse morphism       & Parity          & work \\
$\bowtie$  & composition            & Fidelity        & work \\
$\top$     & evaluation, truth axis & Kinetics        & work \\
$\bot$     & evaluation, falsity axis & Chirality     & work \\
$\boxplus$ & evaluation, information axis & Stoichiometry & work \\
$\Box$     & fixpoint               & Protection      & work \\
\bottomrule
\end{tabular}
\end{center}

A \emph{word} is a finite sequence of marks. Words are read as loops: the last mark
is followed by the first. This is forced rather than chosen, and we return to it in
\Cref{sec:cyclic}.

\subsection{Four values}

The readout takes values in the four-element Belnap lattice \cite{belnap1977,
dunn1976}, whose elements record not only truth and falsity but the two ways a
determination can fail: $B$, both, an overdetermination that holds a conflict; and
$N$, neither, an underdetermination that holds nothing. The four values are used
here in their structural reading:

\begin{itemize}
\item $T$ --- the word closes. Every fork pairs with a fuse, and at least one
      paired region encloses work. Something was split, something was done to it,
      and the arms were rejoined.
\item $B$ --- a fork is left open. The word divides and does not rejoin.
\item $N$ --- forks and fuses pair, but no paired region encloses work. The word
      divides and rejoins having done nothing between.
\item $F$ --- a fuse has no fork. The word rejoins what was never divided, which is
      ill-typed.
\end{itemize}

$N$ and $F$ are distinct failures and must not be collapsed. $F$ is a type error;
$N$ is a well-formed word that performs no transformation.

\subsection{Cyclic pairing}\label{sec:cyclic}

\begin{definition}[Pairing]
\label{def:pairing}
Write $b(i)$ for the number of forks minus the number of fuses among the first $i$
marks of $w$, so $b(0) = 0$ and $b(n) = d$ where $n$ is the length of $w$ and
$d = |w|_{\in} - |w|_{\ni}$. Let $k$ be an index at which $b$ attains its minimum.
Read $w$ as a cycle beginning at position $k$; on each $\in$ push its position onto
a stack, and on each $\ni$ pop, forming a pair whose \emph{region} runs forward
around the cycle from the fork to the fuse. A fuse that finds the stack empty is
\emph{unmatched}; positions left on the stack after the full traversal are the
\emph{unmatched forks}.
\end{definition}

The rotation is not a convenience. A word is a loop, and reading from an arbitrary
cut makes a fuse whose fork lies across the cut appear unmatched when its region is
in fact intact; the value would then depend on where one began. Beginning at a
minimum of $b$ removes every such artefact at once, as \Cref{lem:surplus} shows.

\begin{definition}[Readout]
\label{def:readout}
For a word $w$ with pairing as above:
\begin{enumerate}
\item if $w$ contains no fork and no fuse, the value is $N$;
\item otherwise, if there are more unmatched fuses than unmatched forks, the value
      is $F$;
\item otherwise, if there is an unmatched fork, the value is $B$;
\item otherwise, if some region contains a work mark, the value is $T$;
\item otherwise the value is $N$.
\end{enumerate}
\end{definition}

Pairing must be cyclic. A word is a loop and the cyclic shift is a symmetry of the
reading, so a rotation that cuts through the interior of a region must not change
the value. Pairing over the linearised list makes it change: a fuse landing before
its fork in the cut reads as ill-typed while the region is in fact intact, and the
value becomes a function of where one began reading rather than a property of the
word. Every statement below is about the cyclic pairing.

\section{The readout in closed form}

\Cref{def:readout} is a procedure. The following two results replace it by a
count and a local condition.

\begin{lemma}[Surplus]
\label{lem:surplus}
For any word $w$, with $d = |w|_{\in} - |w|_{\ni}$,
\[
  \#\{\text{unmatched forks}\} = \max(0, d),
  \qquad
  \#\{\text{unmatched fuses}\} = \max(0, -d).
\]
In particular at most one of the two sets is non-empty, and which one is fixed by
the sign of $d$.
\end{lemma}

\begin{proof}
Suppose first $d \ge 0$. Reading from $k$, the balance after $j$ steps is
$b(k + j) - b(k)$ while $k + j \le n$, which is non-negative because $b$ attains
its minimum at $k$; and once the read wraps, writing $j' = k + j - n$, it is
$b(n) - b(k) + b(j') = d + \big(b(j') - b(k)\big) \ge d \ge 0$, again by minimality.
So the running balance never falls below zero, and a fuse never finds the stack
empty: no fuse is unmatched. Each pair consumes one fork with one fuse, so the
forks left on the stack number $|w|_{\in} - |w|_{\ni} = d$.

For $d < 0$, apply the same argument to the word read backwards with the roles of
fork and fuse exchanged. That reading is the mirror of the one above and pairs the
same regions in the opposite order, so it has the same pairs; its surplus is
$-d > 0$ fuses, and no fork is unmatched.
\end{proof}

\Cref{lem:surplus} is the whole content of the $B$/$F$ distinction. It says the two
failures cannot coexist, that the surplus species is determined by the sign of a
difference of counts, and --- since \Cref{def:readout} tests only whether an
unmatched set is non-empty --- that the magnitude of the surplus is invisible to the
readout.

\begin{theorem}[Closed form]
\label{thm:closed}
Let $w$ be a word, let $d = |w|_{\in} - |w|_{\ni}$, and call a paired region
\emph{substantial} if it contains at least one of the seven work marks
$\succ, \prec, \bowtie, \top, \bot, \boxplus, \Box$. Then the readout of $w$ is
\[
  \begin{cases}
    B & d > 0,\\
    F & d < 0,\\
    T & d = 0 \text{ and some region is substantial},\\
    N & d = 0 \text{ and no region is substantial}.
  \end{cases}
\]
\end{theorem}

\begin{proof}
If $d > 0$ then by \Cref{lem:surplus} there is an unmatched fork and no unmatched
fuse, so clause 2 of \Cref{def:readout} does not fire and clause 3 does, giving $B$.
If $d < 0$ there is an unmatched fuse and no unmatched fork, so clause 2 fires,
giving $F$. If $d = 0$ both unmatched sets are empty, and clauses 4 and 5 decide by
exactly the substantiality condition. The case of no forks and no fuses is
$d = 0$ with no regions at all, hence no substantial region, and clause 1 and clause
5 agree on $N$.
\end{proof}

\begin{remark}
The four values are therefore not four independent conditions on a word but a
single integer together with one existential condition on it. In particular the
value cannot be shifted from $B$ to $F$ or back by any rearrangement, since
rearrangement preserves $d$; and it cannot be shifted from $B$ or $F$ to $T$ by any
rearrangement either. Only a change of composition moves a word off $d \neq 0$,
and only an arrangement question remains once $d = 0$.
\end{remark}

\subsection{The enclosure condition is about work, not about length}

Substantiality is not the same as a region being non-empty. Five marks do no work:
the two boundaries, self-reference, and the fork and fuse themselves. A region
containing only these is not substantial, and the word reads $N$.

The distinction is easy to miss because it is invisible on natural sequences. In the
corpus of \Cref{sec:verification}, every non-empty region contains a work mark, so
a weaker reading of \Cref{thm:closed} --- ``some region is non-empty'' --- agrees
with the correct one on every word tested. The two are separated by an explicit
witness:
\[
  \vdash\, \in \odot \ni \,\dashv \quad\text{reads } N,
  \qquad
  \vdash\, \in \bowtie \ni \,\dashv \quad\text{reads } T.
\]
The first region is non-empty and not substantial. Its content is self-reference: a
fork that encloses only $\odot$ has looked at itself and returned. Repeating the
self-reference does not help, and neither does enclosing a boundary; the value stays
$N$ until a work mark appears inside a region.

This is what makes $N$ a coherent value rather than a residue. $N$ is not a fold
that failed. It is a fold that never opened.

\section{The lift from a transcript}

\subsection{Bases}

The four bases carry to the four Belnap values by their pairing behaviour. Guanine
wobble-pairs with both cytosine and uracil, so it carries the overdetermined value;
cytosine pairs only with guanine and carries the true value; adenine carries the
false and uracil the neither. The assignment is a bijection of four elements onto
four.

\subsection{Codons}

The standard genetic code \cite{crick1968, nirenberg1965} takes the sixty-four
codons to twenty amino acids and three stops.

\subsection{Amino acids}

Exactly twelve of the twenty amino acids are promoted, and the promoted set is in
bijection with the twelve marks. The other eight activate no axis and emit nothing;
this is a property of the code, not a gap in the reading. Twenty-three of the
sixty-one sense codons are promoted.

\begin{center}
\begin{tabular}{cllll}
\toprule
Mark & Residue & Axis & Codons & Class \\
\midrule
$\vdash$   & Met & Dimensionality & AUG          & no work \\
$\dashv$   & Trp & Topology       & UGG          & no work \\
$\odot$    & His & Criticality    & CAU, CAC     & no work \\
$\in$      & Asn & Granularity    & AAU, AAC     & no work \\
$\ni$      & Gln & Grammar        & CAA, CAG     & no work \\
\midrule
$\succ$    & Cys & Recognition    & UGU, UGC     & work \\
$\prec$    & Tyr & Parity         & UAU, UAC     & work \\
$\bowtie$  & Phe & Fidelity       & UUU, UUC     & work \\
$\top$     & Ile & Kinetics       & AUU, AUC, AUA & work \\
$\bot$     & Asp & Chirality      & GAU, GAC     & work \\
$\boxplus$ & Lys & Stoichiometry  & AAA, AAG     & work \\
$\Box$     & Glu & Protection     & GAA, GAG     & work \\
\bottomrule
\end{tabular}
\end{center}

A coding sequence is read from its first in-frame AUG and stops at a stop codon;
each promoted codon emits its mark and each unpromoted one emits nothing. The
result is the word of the sequence.

\section{The readout of a coding sequence}

The fork is asparagine and the fuse is glutamine. Their codon families are
\[
  \in \;=\; \mathrm{AAY} \;=\; \{\mathrm{AAU}, \mathrm{AAC}\},
  \qquad
  \ni \;=\; \mathrm{CAR} \;=\; \{\mathrm{CAA}, \mathrm{CAG}\}.
\]
They differ in the first base alone --- adenine forks, cytosine fuses --- with
adenine fixed at the second position. The third base only selects between the two
families sharing the first two, and the full $\{A,C\}A\{U,C,A,G\}$ block splits four
ways:
\[
  \mathrm{AAY} \to \in \text{ fork},\quad
  \mathrm{AAR} \to \boxplus \text{ Lys},\quad
  \mathrm{CAR} \to \ni \text{ fuse},\quad
  \mathrm{CAY} \to \odot \text{ His}.
\]

\begin{tcolorbox}[colback=white,colframe=black!70,arc=1pt,boxrule=0.8pt]
\begin{theorem}[Readout of a coding sequence]
\label{thm:codon}
Let $S$ be an open reading frame, and let $a$ and $g$ be the numbers of
$\mathrm{AAY}$ and $\mathrm{CAR}$ codons it contains. Then the readout of $S$ is
\begin{itemize}
\item $B$ if $a > g$;
\item $F$ if $a < g$;
\item $T$ if $a = g$ and some $\mathrm{AAY}$ codon is separated, along the frame
      read cyclically, from the $\mathrm{CAR}$ codon it pairs with by at least one
      codon of Cys, Tyr, Phe, Ile, Asp, Lys or Glu;
\item $N$ if $a = g$ and no such separation occurs.
\end{itemize}
No other feature of the sequence --- length, composition in the other seventeen
residues, or the order of anything outside a paired region --- affects the value.
\end{theorem}
\end{tcolorbox}

\begin{proof}
Immediate from \Cref{thm:closed} and the tables of Section 3, since the lift sends
$\mathrm{AAY}$ to $\in$, $\mathrm{CAR}$ to $\ni$, the seven listed residues to the
seven work marks, and Met, Trp and His to the three promoted marks that do no work.
\end{proof}

The three no-work promoted residues besides the fork and the fuse are methionine,
tryptophan and histidine. An asparagine and a glutamine separated only by these ---
or by nothing at all --- enclose no work, and the sequence reads $N$.

\subsection{An instance}

The two chains of insulin are the sharpest available illustration, because they
carry the same fork count and differ only in what the fork encloses. The B chain
lifts to
\[
  \bowtie\; \in \ni \;\odot \succ \odot \Box \prec \succ \Box \bowtie \bowtie \prec \boxplus,
\]
in which the asparagine is immediately followed by the glutamine. The region is
empty; the six work marks in the chain all lie outside it; the chain reads $N$. The
A chain, at the same fork count, has regions enclosing three and twelve marks and
reads $T$. Work adjacent to a fork is not work that the fork does.

\section{Verification}\label{sec:verification}

\subsection{Against the corpus}

Sixty-five deposited structures bearing a residue sequence were lifted --- one mark
per resolved $\alpha$-carbon, promoted residues only --- and read out by
\Cref{def:readout}. The distribution is $9$ words at $T$, $22$ at $B$, $13$ at $N$
and $21$ at $F$. The pairing procedure of \Cref{def:pairing} and the closed form of
\Cref{thm:closed} return the same value on all sixty-five, and that value is the
recorded one in every case. The Surplus Lemma was checked on the same corpus in the
form it is stated: on each word exactly one unmatched set is empty and the other has
the size $|d|$.

The corpus separates the surplus claim sharply. Every $B$ word has positive
$d$ and every $F$ word negative, with $d$ ranging over $+1$ to $+36$ on the $B$
side and $-1$ to $-5$ on the $F$ side, and the value constant across each range.
The magnitude is invisible, as \Cref{lem:surplus} requires.

At $d = 0$ the corpus splits nine to thirteen, and every one of the thirteen $N$
words either contains no fork at all or contains a fork adjacent to its fuse. As
noted above, the corpus does not separate substantiality from non-emptiness; that
separation is carried by the explicit witness and not by the data.

\subsection{Against reconstructed transcripts}

For nine of the structures, the residue sequence was read from the deposited
coordinates, back-translated to a coding sequence, and lifted through the codon
route of \Cref{thm:codon}. The value agrees with the structural route in all nine
cases:

\begin{center}
\begin{tabular}{lrrrc}
\toprule
Structure & $a$ & $g$ & $d$ & Value \\
\midrule
Trp-cage                 & 1 & 1 & $0$  & $T$ \\
Villin headpiece subdomain & 2 & 2 & $0$  & $T$ \\
Insulin                  & 3 & 3 & $0$  & $T$ \\
$\beta$-endorphin        & 2 & 2 & $0$  & $T$ \\
A designed binder that closes & 4 & 4 & $0$  & $T$ \\
Insulin B chain          & 1 & 1 & $0$  & $N$ \\
Green fluorescent protein & 13 & 7 & $+6$ & $B$ \\
Ubiquitin                & 2 & 6 & $-4$ & $F$ \\
A designed binder that does not & 2 & 6 & $-4$ & $F$ \\
\bottomrule
\end{tabular}
\end{center}

Trp-cage is the cleanest case. The transcript yields the word
$\vdash \in \prec \top \ni \dashv \boxplus \bot$ once; the deposited structure, an
ensemble of thirty-eight models, yields the same motif thirty-eight times. The
sequence carries the reading, and the structure file carries it once per deposited
copy.

\section{Position}

\subsection{What is settled}

The four-valued readout is a function of two counts and one existential condition,
and \Cref{thm:closed} states it in closed form with a proof from the pairing
definition rather than from the corpus. Carried through the genetic code, the
readout of an open reading frame is determined by its $\mathrm{AAY}$ and
$\mathrm{CAR}$ codons and by which of seven other codon families separate a paired
asparagine from its glutamine. The two failure modes are decided by composition
alone and are invariant under every rearrangement of the sequence; only the
distinction between closure and the null reading is a question of arrangement.

Three consequences are worth stating separately. First, the value cannot be
repaired by insertion or by permutation when $d \neq 0$: permutation preserves $d$,
and insertion that fixes $d$ changes the composition, which is a rebuild and not a
repair. Second, the value is computed from the promoted subsequence alone, so it is
fixed at translation and nothing downstream of translation contributes to it.
Third, $N$ is not a degenerate $T$ but a separate structural fact, and the witness
$\vdash \in \odot \ni \dashv$ shows that its boundary with $T$ is drawn at work and
not at length.

\subsection{What is open}

The readout says which of four structural verdicts a sequence receives. It does not
say what those verdicts predict about the folded molecule, and this paper claims no
such correspondence: the agreement reported in \Cref{sec:verification} is between
two routes to the same readout, not between the readout and any thermodynamic or
kinetic observable. Whether the classes align with foldedness, with stability, or
with any measured quantity is an empirical question that the closed form now makes
cheap to ask, since the classification of an arbitrary reading frame costs a single
pass counting two codon families.

Two structural questions also remain. The condition at $d = 0$ is existential ---
one substantial region suffices, and a word with many empty regions and one
substantial region reads $T$ --- so the readout does not distinguish a sequence with
a single working fork from one in which every fork works. A refinement counting
substantial regions rather than testing for one would separate these, and its
relation to the four values is not settled here. And the alphabet's cyclic reading
means the first and last marks of a frame are adjacent; what that adjacency
corresponds to in a linear polypeptide is a question this paper leaves open, having
used the cyclic pairing throughout because the shift symmetry demands it.

\begin{thebibliography}{9}

\bibitem{belnap1977} N.~D.~Belnap.
\emph{A useful four-valued logic}. In J.~M.~Dunn and G.~Epstein, editors, Modern
Uses of Multiple-Valued Logic, pages 5--37. Reidel, Dordrecht, 1977.

\bibitem{dunn1976} J.~M.~Dunn.
\emph{Intuitive semantics for first-degree entailments and coupled trees}.
Philosophical Studies, 29(3):149--168, 1976.

\bibitem{crick1968} F.~H.~C.~Crick.
\emph{The origin of the genetic code}. Journal of Molecular Biology,
38(3):367--379, 1968.

\bibitem{nirenberg1965} M.~Nirenberg, P.~Leder, M.~Bernfield, R.~Brimacombe,
J.~Trupin, F.~Rottman, and C.~O'Neal.
\emph{RNA codewords and protein synthesis, VII. On the general nature of the RNA
code}. Proceedings of the National Academy of Sciences, 53(5):1161--1168, 1965.

\bibitem{anfinsen1973} C.~B.~Anfinsen.
\emph{Principles that govern the folding of protein chains}. Science,
181(4096):223--230, 1973.

\bibitem{levinthal1968} C.~Levinthal.
\emph{Are there pathways for protein folding?} Journal de Chimie Physique,
65:44--45, 1968.

\bibitem{neidigh2002} J.~W.~Neidigh, R.~M.~Fesinmeyer, and N.~H.~Andersen.
\emph{Designing a 20-residue protein}. Nature Structural Biology, 9(6):425--430,
2002.

\bibitem{mcknight1997} C.~J.~McKnight, D.~S.~Doering, P.~T.~Matsudaira, and
P.~S.~Kim.
\emph{A thermostable 35-residue subdomain within villin headpiece}. Journal of
Molecular Biology, 260(2):126--134, 1996.

\bibitem{vijaykumar1987} S.~Vijay-Kumar, C.~E.~Bugg, and W.~J.~Cook.
\emph{Structure of ubiquitin refined at 1.8~\AA{} resolution}. Journal of Molecular
Biology, 194(3):531--544, 1987.

\bibitem{ormo1996} M.~Orm\"o, A.~B.~Cubitt, K.~Kallio, L.~A.~Gross, R.~Y.~Tsien,
and S.~J.~Remington.
\emph{Crystal structure of the Aequorea victoria green fluorescent protein}.
Science, 273(5280):1392--1395, 1996.

\bibitem{priest1979} G.~Priest.
\emph{The logic of paradox}. Journal of Philosophical Logic, 8(1):219--241, 1979.

\end{thebibliography}

\end{document}
